\documentclass[../../main.tex]{subfiles}

\begin{document}

\section{子空间拓扑}\label{section:ch02-子空间拓扑}

\begin{definition}[子空间拓扑]\label{definition:ch02-子空间拓扑}
设 $X$ 是一个拓扑空间，其拓扑为 $\mathcal{T}$。若 $Y$ 是 $X$ 的一个子集，则族
\begin{align*}
\mathcal{T}_{Y} = \{Y \cap U \mid U \in \mathcal{T}\}
\end{align*}
是 $Y$ 的一个拓扑，称为子空间拓扑（subspace topology）。具有这种拓扑的 $Y$ 称为 $X$ 的一个子空间（subspace），其开集由 $X$ 中的开集与 $Y$ 的交组成。
\end{definition}

容易看出，$\mathcal{T}_{Y}$ 是一个拓扑。因为
\begin{align*}
\varnothing = Y \cap \varnothing, \quad Y = Y \cap X,
\end{align*}
所以 $\varnothing$ 和 $Y$ 是 $\mathcal{T}_{Y}$ 的元素（其中 $\varnothing$ 和 $X$ 是 $\mathcal{T}$ 的元素）。由等式
\begin{align*}
(U_{1} \cap Y) \cap \cdots \cap (U_{n} \cap Y) = (U_{1} \cap \cdots \cap U_{n}) \cap Y,
\end{align*}
\begin{align*}
\bigcup_{\alpha \in J} (U_{\alpha} \cap Y) = (\bigcup_{\alpha \in J} U_{\alpha}) \cap Y
\end{align*}
可以得到：$\mathcal{T}_{Y}$ 对于有限交、任意并的运算是封闭的。

\begin{lemma}[子空间拓扑的基]\label{lemma:ch02-子空间基}
若 $\mathcal{B}$ 是 $X$ 的拓扑的一个基，则族
\begin{align*}
\mathcal{B}_{Y} = \{B \cap Y \mid B \in \mathcal{B}\}
\end{align*}
是 $Y$ 上子空间拓扑的一个基。
\end{lemma}

\begin{proof}
给定 $X$ 的一个开集 $U$ 及 $y \in U \cap Y$，我们能在 $\mathcal{B}$ 中选取一个元素 $B$，使得 $y \in B \subset U$。因此，$y \in B \cap Y \subset U \cap Y$。根据\cref{lemma:ch02-开集族为基}可见，$\mathcal{B}_{Y}$ 是 $Y$ 的子空间拓扑的一个基。
\end{proof}

当我们处理空间 $X$ 和子空间 $Y$ 的时候，使用“开集”这个词必须特别小心。究竟它指的是 $Y$ 的拓扑的一个元素，还是 $X$ 的拓扑的一个元素？我们作出以下规定，若 $Y$ 是 $X$ 的一个子空间，集合 $U$ 属于 $Y$ 的拓扑，则称 $U$ 是 $Y$ 中的一个开集（或者说是相对于 $Y$ 而言的一个开集），特别地，这意味着 $U$ 是 $Y$ 的一个子集。如果 $U$ 属于 $X$ 的拓扑，则称 $U$ 是 $X$ 中的一个开集。

在一种特殊情况下，$Y$ 中的任何一个开集也是 $X$ 中的开集：

\begin{lemma}[子空间开集为全空间开集的条件]\label{lemma:ch02-子空间开集条件}
设 $Y$ 是 $X$ 的一个子空间，若 $U$ 是 $Y$ 的一个开集并且 $Y$ 是 $X$ 的一个开集，则 $U$ 是 $X$ 的一个开集。
\end{lemma}

\begin{proof}
因为 $U$ 是 $Y$ 的一个开集，则存在 $X$ 的某一个开集 $V$，使得 $U = Y \cap V$。又因为 $Y$ 与 $V$ 都是 $X$ 的开集，所以 $Y \cap V$ 也是 $X$ 的开集。
\end{proof}

以下我们讨论子空间拓扑与序拓扑和积拓扑之间的关系。我们将看到：有关积拓扑的相应结论与人们的期望相符，而关于序拓扑的相关结论则不然。

\begin{theorem}[积拓扑与子空间拓扑一致]\label{theorem:ch02-积子空间一致}
若 $A$ 是 $X$ 的一个子空间，$B$ 是 $Y$ 的一个子空间，则 $A \times B$ 的积拓扑与它从 $X \times Y$ 继承的子空间拓扑是同一个拓扑。
\end{theorem}

\begin{proof}
集合 $U \times V$ 是 $X \times Y$ 的一个基元素，其中 $U$ 是 $X$ 的一个开集并且 $V$ 是 $Y$ 的一个开集。于是 $(U \times V) \cap (A \times B)$ 是 $A \times B$ 的子空间拓扑的一个基元素。注意：
\begin{align*}
(U \times V) \cap (A \times B) = (U \cap A) \times (V \cap B).
\end{align*}
因为 $U \cap A$ 和 $V \cap B$ 分别是 $A$ 和 $B$ 上子空间拓扑的开集，所以集合 $(U \cap A) \times (V \cap B)$ 是 $A \times B$ 的积拓扑的一个基元素。这样我们得到：作为 $A \times B$ 上子空间拓扑的基与作为 $A \times B$ 上积拓扑的基是一样的，因此这两个拓扑相同。
\end{proof}

以下设 $X$ 为具有序拓扑的全序集，$Y$ 是 $X$ 的一个子集。将 $X$ 上的序关系限制在 $Y$ 上使 $Y$ 成为全序集。然而，由此得到的 $Y$ 上的序拓扑与 $Y$ 从 $X$ 继承的子空间拓扑未必是同一拓扑。下面给出一个二者相同的例子和两个二者不同的例子。

\begin{example}\label{example:ch02-区间子空间}
就子空间拓扑而言，研究实直线 $\mathbb{R}$ 的子集 $Y = [0, 1]$。子空间拓扑的基是所有形如 $(a, b) \cap Y$ 的集合所成的族，其中 $(a, b)$ 是 $\mathbb{R}$ 的一个开区间。这种集合是下面几个类型的一种：
\begin{align*}
(a, b) \cap Y =
\begin{cases}
(a, b) & \text{如果 } a \text{ 和 } b \text{ 都属于 } Y, \\
[0, b) & \text{若仅 } b \text{ 属于 } Y, \\
(a, 1] & \text{若仅 } a \text{ 属于 } Y, \\
Y \text{ 或者 } \varnothing & \text{如果 } a \text{ 和 } b \text{ 都不属于 } Y.
\end{cases}
\end{align*}
根据定义易见，上述每一个集合都是 $Y$ 的开集。但是，第二、三两种类型的集合不是空间 $\mathbb{R}$ 的开集。

注意，上述形式的集合组成了 $Y$ 上序拓扑的一个基。从而对于集合 $Y = [0, 1]$ 的情形，它的子空间拓扑（作为 $\mathbb{R}$ 的一个子空间）与序拓扑是一样的。
\end{example}

\begin{example}\label{example:ch02-子空间不同1}
设 $Y$ 是 $\mathbb{R}$ 的子集 $[0, 1) \cup \{2\}$。对于 $Y$ 的子空间拓扑而言，单点集 $\{2\}$ 是开集。这是由于它是开集 $\left(\frac{3}{2}, \frac{5}{2}\right)$ 与 $Y$ 的交的缘故。但是对于 $Y$ 的序拓扑而言，集合 $\{2\}$ 不是开集。对于 $Y$ 序拓扑的任何包含 2 的基元素，形如
\begin{align*}
\{x \mid x \in Y \quad \text{和} \quad a < x \leqslant 2\},
\end{align*}
其中 $a$ 为 $Y$ 中某一点。这种集合必定包含 $Y$ 中小于 2 的点。
\end{example}

\par
\begin{figure}[htbp]
\centering
\includegraphics[width=0.6\textwidth]{ch02-fig16-1}
\caption{子空间拓扑与序拓扑的比较}\label{figure:ch02-子空间vs序}
\end{figure}

\begin{example}\label{example:ch02-子空间不同2}
设 $I = [0, 1]$。由 $\mathbb{R} \times \mathbb{R}$ 上字典序的限制得到 $I \times I$ 上的字典序。然而 $I \times I$ 上的序拓扑与它从 $\mathbb{R} \times \mathbb{R}$ 上继承的子空间拓扑是不同的拓扑！例如，$\{1/2\} \times (1/2, 1]$ 是 $I \times I$ 的子空间拓扑的开集，但是，易见它不是序拓扑中的开集。参见\cref{figure:ch02-子空间vs序}。


具有字典序拓扑的拓扑空间 $I \times I$ 称为有序矩形（ordered square），记作 $I_{o}^{2}$。
\end{example}

对于全序集 $X$ 上的一个区间或一条射线而言，将不会出现\cref{example:ch02-子空间不同1}和\cref{example:ch02-子空间不同2}所示的不正常情况。见下面的定理。

对于全序集 $X$，它的一个子集 $Y$ 称为在 $X$ 中是凸的（convex），若对于 $Y$ 中的每一个点的偶对 $a < b$，$X$ 中的整个区间 $(a, b)$ 包含于 $Y$。注意，$X$ 中的区间和射线都是 $X$ 中的凸子集。

\begin{theorem}[凸子集上序拓扑与子空间拓扑一致]\label{theorem:ch02-凸序一致}
设 $X$ 是具有序拓扑的一个全序集，并且 $Y$ 是 $X$ 的一个凸子集，那么 $Y$ 的序拓扑与它从 $X$ 继承的子空间拓扑是同一个拓扑。
\end{theorem}

\begin{proof}
设 $(a, +\infty)$ 是 $X$ 的一条射线，考虑它与 $Y$ 的交。若 $a \in Y$，则
\begin{align*}
(a, +\infty) \cap Y = \{x \mid x \in Y \quad \text{和} \quad x > a\}
\end{align*}
为全序集 $Y$ 中的一条开射线。若 $a \notin Y$，由于 $Y$ 是一个凸子集，则或者 $a$ 是 $Y$ 的一个下界，或者 $a$ 是 $Y$ 的一个上界。对于前者，$(a, +\infty) \cap Y$ 等于空间 $Y$。对于后者，$(a, +\infty) \cap Y$ 为空集。

类似的讨论说明：射线 $(-\infty, a)$ 与 $Y$ 的交或者为 $Y$ 中的一条开射线，或者为 $Y$ 自身，或者为空集。由于所有形如 $(a, +\infty) \cap Y$ 和 $(-\infty, a) \cap Y$ 的集合已构成了 $Y$ 的子空间拓扑的一个子基，并且它们中的每一个都是序拓扑的开集，从而序拓扑包含了子空间拓扑。

为了证明反过来的包含关系，注意 $Y$ 的每一条开射线都是 $X$ 中的一条开射线与 $Y$ 的交，因而它是 $Y$ 的子空间拓扑的开集。由于 $Y$ 的开射线构成了 $Y$ 的序拓扑的一个子基，从而序拓扑包含于子空间拓扑。
\end{proof}

为了避免混淆，我们作这样的约定：当讨论全序集 $X$ 的子集 $Y$ 时，除非有特别的说明，总是假设 $Y$ 取定的是子空间拓扑。若 $Y$ 是 $X$ 中的凸子集，子空间拓扑与序拓扑相同；否则，它们可能不同。

\subsection*{习题}

\begin{exercise}\label{exercise:ch02-子空间拓扑-1}
证明：若 $Y$ 是 $X$ 的一个子空间并且 $A$ 是 $Y$ 的一个子集，则 $A$ 从 $Y$ 继承的子空间拓扑与 $A$ 从 $X$ 继承的子空间拓扑是同一个拓扑。
\end{exercise}

\begin{exercise}\label{exercise:ch02-子空间拓扑-2}
若 $\mathcal{T}$ 和 $\mathcal{T}'$ 是 $X$ 上的两个拓扑，并且 $\mathcal{T}'$ 严格细于 $\mathcal{T}$，那么对于 $X$ 的子集 $Y$ 上的相应的子空间拓扑有些什么结论？
\end{exercise}

\begin{exercise}\label{exercise:ch02-子空间拓扑-3}
考虑作为 $\mathbb{R}$ 子空间的集合 $Y = [-1, 1]$，下面集合中哪些是 $Y$ 的开集？哪些是 $\mathbb{R}$ 的开集？
\begin{align*}
A = \{x \mid \frac{1}{2} < |x| < 1\},
\end{align*}
\begin{align*}
B = \{x \mid \frac{1}{2} < |x| \leqslant 1\},
\end{align*}
\begin{align*}
C = \{x \mid \frac{1}{2} \leqslant |x| < 1\},
\end{align*}
\begin{align*}
D = \{x \mid \frac{1}{2} \leqslant |x| \leqslant 1\},
\end{align*}
\begin{align*}
E = \{x \mid 0 < |x| < 1 \text{ 和 } 1/x \notin \mathbb{Z}_{+}\}.
\end{align*}
\end{exercise}

\begin{exercise}\label{exercise:ch02-子空间拓扑-4}
映射 $f: X \rightarrow Y$ 称为一个开映射（open map），如果对于 $X$ 的每一个开集 $U$，集合 $f(U)$ 是 $Y$ 中的一个开集。证明 $\pi_{1}: X \times Y \rightarrow X$ 及 $\pi_{2}: X \times Y \rightarrow Y$ 都是开映射。
\end{exercise}

\begin{exercise}\label{exercise:ch02-子空间拓扑-5}
设 $X$ 和 $X'$ 分别表示具有拓扑 $\mathcal{T}$ 和 $\mathcal{T}'$ 的同一个集合。$Y$ 和 $Y'$ 表示分别具有拓扑 $\mathcal{U}$ 和 $\mathcal{U}'$ 的同一个集合。

(a) 证明：若 $\mathcal{T}' \supset \mathcal{T}$ 并且 $\mathcal{U}' \supset \mathcal{U}$，则 $X' \times Y'$ 上的积拓扑细于 $X \times Y$ 上的积拓扑。

(b) (a) 的逆是否成立？验证你的结论。
\end{exercise}

\begin{exercise}\label{exercise:ch02-子空间拓扑-6}
证明可数族
\begin{align*}
\{(a, b) \times (c, d) \mid a < b, c < d, \ \text{其中 } a, b, c, d \text{ 都是有理数}\}
\end{align*}
是 $\mathbb{R}^{2}$ 的一个基。
\end{exercise}

\begin{exercise}\label{exercise:ch02-子空间拓扑-7}
设 $X$ 为全序集。试问：$Y$ 为 $X$ 的一个凸子集，是否意味着 $Y$ 是 $X$ 的一个区间或射线？
\end{exercise}

\begin{exercise}\label{exercise:ch02-子空间拓扑-8}
设 $L$ 为平面上的一条直线，试描述 $L$ 分别从 $\mathbb{R}_{\ell} \times \mathbb{R}$ 和 $\mathbb{R}_{\ell} \times \mathbb{R}_{\ell}$ 继承的子空间拓扑。这两种拓扑都是我们所熟悉的拓扑。
\end{exercise}

\begin{exercise}\label{exercise:ch02-子空间拓扑-9}
证明 $\mathbb{R} \times \mathbb{R}$ 上的字典序拓扑与积拓扑 $\mathbb{R}_{d} \times \mathbb{R}$ 是同一拓扑，这里 $\mathbb{R}_{d}$ 表示 $\mathbb{R}$ 上的离散拓扑。并将这个拓扑与 $\mathbb{R}^{2}$ 上的标准拓扑进行比较。
\end{exercise}

\begin{exercise}\label{exercise:ch02-子空间拓扑-10}
设 $I = [0, 1]$。试比较 $I \times I$ 的积拓扑，$I \times I$ 上的字典序拓扑，$I \times I$ 从 $\mathbb{R} \times \mathbb{R}$ 的字典序拓扑继承而来的子空间拓扑。
\end{exercise}

\end{document}